% pre-warm dummy — compile ตอน Docker build เพื่อฝัง TeX package cache ใน image
% ครอบคลุม package ที่ report ใช้บ่อย (compile prod แรกจะไม่ต้องดึงเน็ต)
\documentclass[a4paper,12pt]{article}
\usepackage{fontspec}
\usepackage{geometry}
\usepackage{longtable}
\usepackage{array}
\usepackage{booktabs}
\usepackage{xcolor}
\usepackage{graphicx}
% เอกสารราชการ/แบบฟอร์ม (2026-09-22): จัดตำแหน่งแน่นอน · พื้นหลังแบบฟอร์ม · ตาราง · ขีดเส้นใต้
\usepackage{tabularx}
\usepackage{multirow}
\usepackage[normalem]{ulem}
\usepackage{fancyhdr}
\usepackage{eso-pic}
\usepackage[absolute,overlay]{textpos}
\usepackage{tikz}
\geometry{margin=2cm}
\setmainfont{Sarabun}
\begin{document}
\section*{pre-warm}
ทดสอบภาษาไทย\hskip0pt plus0.5pt minus0.3pt\relax และ glue break — English 123 ๑๒๓
\begin{longtable}{|l|r|}
\hline รายการ & จำนวน \\ \hline
A & 1 \\ \hline
\end{longtable}
\textcolor{red}{color test}
\begin{tikzpicture}[remember picture,overlay]\node at (current page.center) {tikz};\end{tikzpicture}
\begin{textblock*}{3cm}(2cm,2cm)textpos\end{textblock*}
\begin{tabularx}{\linewidth}{|X|r|}\hline \multirow{1}{*}{x} & \uline{1} \\ \hline\end{tabularx}
% ขนาดตัวอักษรเอกสารราชการ (16pt/29pt) — ให้ฟอนต์ math ขนาดนั้นอยู่ใน cache ด้วย
{\fontsize{16pt}{18pt}\selectfont ขนาด 16 $x^2$\\}{\fontsize{29pt}{32pt}\selectfont ขนาด 29 $x$\\}
$\left\{\rule{0pt}{1cm}\right.\sum$ % วงเล็บใหญ่/สัญลักษณ์ math (cmex10)
\end{document}
